%------------------------------------------
% Trabalho Relacionados
%------------------------------------------
\section{Trabalhos Relacionados}
\label{sec:trabalhos-re}


\todo[inline]{Nesta parte faz-se a revisão da literatura e do estado da prática sobre o assunto, analisando criticamente os resultados de estudos e soluções desenvolvidos por outros autores, cujas obras citadas e consultadas devem constar nas referências.}

\textcolor{red}{A função desta seção é demonstrar à banca que você conhece o estado da arte e o estado da prática do seu tema e que a lacuna identificada é real. O que faltar de conhecimento ou de solução é lacuna e, portanto, a oportunidade que você explora como problema de pesquisa. No mestrado, espera-se uma revisão estruturada e reprodutível: descreva brevemente o protocolo de busca adotado, isto é, as bases consultadas (por exemplo, IEEE Xplore, ACM Digital Library, Scopus, Web of Science e SBC OpenLib), a \textit{string} de busca, o período considerado e os critérios de inclusão e exclusão. Quando pertinente, conduza um mapeamento sistemático ou uma revisão sistemática da literatura. Priorize artigos de periódicos e conferências relevantes da área, preferencialmente dos últimos cinco anos, sem desconsiderar trabalhos seminais. O Google Acadêmico pode ser usado como fonte complementar.}

\textcolor{red}{As pesquisas que você irá buscar para compor esta seção são aquelas semelhantes à sua, ou seja, que investigam o mesmo tema e que atacam problemas de pesquisa nos quais você está interessado em estudar. Por se tratar de um Mestrado Profissional em Computação Aplicada, considere também o estado da prática: ferramentas, produtos e serviços existentes no mercado ou no setor público, soluções de código aberto, patentes e relatórios técnicos que tratem do mesmo problema no domínio de aplicação. Isso permite demonstrar por que as soluções existentes não atendem ao contexto real em que a sua pesquisa se insere.}

\textcolor{red}{A descrição de um estudo deve ocupar, idealmente, apenas um parágrafo. Como regra geral, um parágrafo pode ter até 20 linhas, mas tente trabalhar com o máximo de 12 ou 16. No máximo, a sua descrição de um estudo deve ocupar dois parágrafos, não mais que isso.}

\textcolor{red}{As descrições de estudos devem ser conectadas entre si por meio do texto. Ou seja, você pode afirmar que Santos (2022) investigou X e concluiu Y, sugerindo como trabalho futuro Z. Aí, no parágrafo seguinte, você pode dizer que Silva (2022), baseado em Santos (2022), conduziu estudo para investigar Z. Note, portanto, que os estudos foram, de algum modo, conectados no seu texto. Agrupe os trabalhos por abordagem, técnica ou subproblema, em vez de apresentá-los como uma lista isolada de resumos.}

\textcolor{red}{Como regra geral, a descrição de um estudo deve seguir o seguinte formato, que você pode adaptar: “Silva e Santos (2022) conduziram estudo cujo objetivo foi {Preencher com o verbo no infinitivo mais complemento adotado por Silva e Santos}. Para atingir esse objetivo, adotaram o seguinte método: {Preencher com uma breve descrição do método do estudo}. Os pesquisadores verificaram que {Preencher com os principais dados}. A partir desses dados, concluíram que {Preencher com a resposta à principal pergunta de pesquisa}. Esse estudo teve como principais limitações: {Preencher com uma breve síntese}. Em relação ao presente trabalho, {Preencher com a semelhança e a diferença em relação à sua proposta}”.}

\textcolor{red}{O parágrafo final desta seção deve explicitar a correlação entre as pesquisas e soluções descritas e, então, apresentar de modo inequívoco a lacuna, científica e/ou prática, que justifica o problema de pesquisa e a solução proposta apresentada na seção seguinte.}

\todo[inline]{Deve-se apresentar uma visão geral e comparativa dos trabalhos apresentados, indicando como cada trabalho pode contribuir com o seu. Sugere-se apresentar uma tabela comparativa entre os trabalhos, com critérios relevantes para o problema e o contexto de aplicação, incluindo uma linha para a sua proposta (``Este trabalho''), de modo a evidenciar o seu diferencial. Abaixo segue um exemplo:}

A seguir analisamos as técnicas (ver Tabela~\ref{Table:TechByPapers}) 
adotadas para a geração das invariantes de programa nos trabalhos apresentados anteriormente. Vale ressaltar que as técnicas identificadas foram:

\begin{itemize}
    \item \textbf{Análise de Predicados}. Analisando os métodos baseados na análise de predicado, podemos observar que a análise de predicado fornece um apoio significativo na análise do programa para inferir propriedades sobre o comportamento do programa;
    %
    \item \textbf{Interpretação Abstrata}. É uma teoria da aproximação da semântica de linguagens de programação cuja aplicação principal é a análise estática;
    %
    \item \textbf{Lógica de Mill}. A lógica de Mill é adotada para caracterizar laços por meio de uma função que define o seu espaço de estados.
    %
    \item \textbf{\textit{Templates}}. Analisando as publicações, notamos que a adoção de \textit{templates} é importante para fornecer um suporte para guiar a inferência de invariantes de programas.
\end{itemize}{}

%\multicolumn{number cols}{align}{text} % align: l,c,r
%\multirow{number rows}{width}{text}
\begin{table}[htbp]  
  \centering  
  \scalefont{0.9}  
  \begin{tabular}{|c|c|c|c|c|}
    \hline
    % header
    % row 1 and 2
    \multicolumn{1}{|c|}{\multirow{2}{*}{\textbf{Artigos}}} & 
    % row 1
    \multicolumn{4}{c|}{\textbf{Técnicas}} \bigstrut\\
    \cline{2-5} &
    % row 2
    \textbf{Predicados} & \textbf{Int. Abstrata} & \textbf{Log. Mill} & \textbf{\textit{Templates}} \bigstrut\\
    \hline
    % body table - row 3
    Trabalho 1 & X &  &  &  \bigstrut\\
    \hline
    % body table - row 4
    Trabalho 2 &  & X & X &  \bigstrut\\
    \hline
    % body table - row 5
    Trabalho 3 & X & X &  & X \bigstrut\\
    \hline
    % body table - row 6
    Trabalho 4 & X & X &  & X \bigstrut\\
    \hline
    % body table - row 7
    \textbf{Este trabalho} & X & X & X & X \bigstrut\\
    \hline
  \end{tabular}  
  \caption{Classificação dos artigos por técnicas.}
  \label{Table:TechByPapers}
\end{table}%
